% Part: computability
% Chapter: computability-theory
% Section: prop-reduce

\documentclass[../../../include/open-logic-section]{subfiles}

\begin{document}

\olfileid{cmp}{thy}{ppr}
\olsection{হ্রাসযোগ্যতার বৈশিষ্ট্য}

$A$, $B$-তে হ্রাস পেলে আমরা $A \leq_m B$ লিখি। এই সংকেতটি বোঝায় যে
$A \leq_m B$ হলে $A$, $B$-এর “চেয়ে কঠিন নয়”, $B$, $A$-এর “সমান কঠিন
বা আরও কঠিন”, এবং সংখ্যার সাধারণ~$\le$-এর মতো $\le_m$ সেটগুলিকে (বা
নির্ণয়-সমস্যাগুলিকে) তাদের জটিলতা অনুসারে বিন্যস্ত করে। নিচের দুটি
প্রস্তাব এই স্বজ্ঞাকে সমর্থন করে। প্রথমটি বলে, $\le_m$ সঞ্চারী।

\begin{prop}
\ollabel{prop:trans-red}
$A \leq_m B$ এবং $B \leq_m C$ হলে $A \leq_m C$।
\end{prop}

\begin{proof}
$A$ থেকে~$B$-তে একটি হ্রাসকরণকে $B$ থেকে~$C$-তে একটি হ্রাসকরণের সঙ্গে
সংযোজিত করলে $A$ থেকে~$C$-তে একটি হ্রাসকরণ পাওয়া যায়।
\end{proof}

\begin{prob}
দেখাও, $f$ ও~$g$ যথাক্রমে $A$ থেকে~$B$ এবং $B$ থেকে~$C$-তে বহু-এক
হ্রাসকরণ হলে $\comp{f}{g}$, $A$ থেকে~$C$-তে একটি বহু-এক হ্রাসকরণ। এভাবে
\olref{prop:trans-red} প্রমাণ করো।
\end{prob}

\begin{prop}
\ollabel{prop:reduce}
$A$ ও $B$ যেকোনো সেট হোক, এবং ধরা যাক $A \leq_m B$।
\begin{enumerate}
\item $B$ গণনসাধ্যভাবে তালিকায়নযোগ্য হলে $A$-ও তাই।
\item $B$ গণনসাধ্য হলে $A$-ও তাই।
\end{enumerate}
\end{prop}

\begin{proof}
$f$-কে $A$ থেকে~$B$-তে একটি বহু-এক হ্রাসকরণ ধরি। প্রথম দাবির জন্য শুধু
যাচাই করো যে $B$ যদি কোনো আংশিক অপেক্ষক~$g$-এর সংজ্ঞাক্ষেত্র হয়, তবে
$A$ হলো~$\comp{f}{g}$-এর সংজ্ঞাক্ষেত্র:
\begin{align*}
x \in A & \text{ যদি এবং কেবল যদি } f(x) \in B \\
& \text{ যদি এবং কেবল যদি }  g(f(x)) \fdefined.
\end{align*}

দ্বিতীয় দাবির জন্য মনে করো, $B$ গণনসাধ্য হলে $B$ ও~$\Complement{B}$
গণনসাধ্যভাবে তালিকায়নযোগ্য (\olref[cmp]{thm:ce-comp})। সহজেই যাচাই করা
যায় যে $f$, $\Complement{A}$ থেকে~$\Complement{B}$-তেও একটি বহু-এক
হ্রাসকরণ। তাই এই প্রমাণের প্রথম অংশ অনুসারে $A$ ও~$\Complement{A}$
!!{computably
enumerable}। অতএব \olref[cmp]{thm:ce-comp} অনুসারে $A$-ও
গণনসাধ্য। (বিকল্পভাবে যাচাই করা যায়, $\Char{A} =
\comp{f}{\Char{B}}$; ফলে $\Char{B}$ গণনসাধ্য হলে~$\Char{A}$-ও তাই।)
\end{proof}

\begin{prob}
ধরা যাক, $f$, $A$ থেকে~$B$-তে একটি বহু-এক হ্রাসকরণ। দেখাও, $f$,
$\Complement{A}$ থেকে~$\Complement{B}$-তেও একটি বহু-এক হ্রাসকরণ।
\end{prob}

\begin{prob}
দেখাও, $f\colon \Nat \to \Nat$ একটি বহু-এক হ্রাসকরণ হলে
$\Char{A} = \comp{f}{\Char{B}}$।
% BN-SRC-153 TARGET-CORRECTION-BEGIN
\par\smallskip\noindent\textnormal{\small\textbf{উৎস-সংশোধন (BN-SRC-153):} হিমায়িত ইংরেজি সমস্যায় $f\colon A\to B$ লেখা হয়েছে, কিন্তু বহু-এক হ্রাসকরণের সংজ্ঞায় $f$-এর ধরন $\Nat\to\Nat$; পূরক ও চরিত্রসূচক অপেক্ষকের সমীকরণটিও সকল স্বাভাবিক নিবেশে এই পূর্ণ ধরন চায়। বাংলা সমস্যায় সংজ্ঞাসঙ্গত ধরনটি রাখা হয়েছে।} % BN-SRC-153 TARGET-CORRECTION-END
\end{prob}

\begin{digress}
\emph{টুরিং হ্রাসযোগ্যতা} নামের হ্রাসযোগ্যতার একটি আরও সাধারণ ধারণা অন্য
প্রসঙ্গে, বিশেষ করে অনির্ণেয়তা প্রমাণে, উপযোগী। লক্ষ করো,
\olref[cmp]{cor:comp-k} অনুসারে~$K_0$-এর পূরক $K_0$-তে হ্রাসযোগ্য নয়,
কারণ সেটি গণনসাধ্যভাবে তালিকায়নযোগ্য নয়। কিন্তু স্বজ্ঞাগতভাবে,
$K_0$-সম্পর্কিত প্রশ্নগুলির উত্তর জানা থাকলে তার পূরক-সম্পর্কিত
প্রশ্নগুলির উত্তরও জানা যেত। কোনো গণনসাধ্য প্রণালী~$B$-সম্পর্কিত প্রশ্ন
করতে পারে এবং সেই প্রণালী ব্যবহার করে~$A$-সম্পর্কিত প্রশ্নের উত্তর
নির্ধারণ করা গেলে $A$-কে~$B$-তে টুরিং হ্রাসযোগ্য বলা হয়। এটি বহু-এক
হ্রাসযোগ্যতার চেয়ে উদার, যেখানে (1)~$B$ সম্পর্কে কেবল একটি প্রশ্ন করা
যায় এবং (2)~“হ্যাঁ” উত্তরকে $A$-সম্পর্কিত প্রশ্নের “হ্যাঁ” উত্তরে এবং
একইভাবে “না” উত্তরকে “না” উত্তরে রূপান্তর করতে হয়। তবু $A$, $B$-তে
টুরিং হ্রাসযোগ্য এবং $B$ গণনসাধ্য হলে $A$-ও গণনসাধ্য (যদিও আমরা দেখেছি,
গণনসাধ্যভাবে তালিকায়নযোগ্যতার জন্য অনুরূপ বক্তব্যটি সত্য নয়)।

আমরা হ্রাসযোগ্যতার যেসব ধারণা আলোচনা করেছি, সেগুলি নিয়ে ভেবে তাদের
পার্থক্যগুলি বোঝা উচিত। তবে এই অধ্যায়ে আমরা কেবল বহু-এক হ্রাসযোগ্যতা
নিয়েই কাজ করব। প্রসঙ্গত, আগের অনুচ্ছেদে আলোচিত দুই ধরনের হ্রাসযোগ্যতারই
গণনামূলক জটিলতায় অনুরূপ ধারণা আছে, যেখানে অতিরিক্ত শর্ত থাকে যে টুরিং
যন্ত্রগুলি বহুপদী সময়ে চলবে: বহু-এক হ্রাসযোগ্যতার জটিলতা-সংস্করণটি
\emph{কার্প হ্রাসযোগ্যতা}, আর টুরিং হ্রাসযোগ্যতার জটিলতা-সংস্করণটি
\emph{কুক হ্রাসযোগ্যতা} নামে পরিচিত।
\end{digress}

\end{document}
